% The bubble-and-arrow language the four l04_dac coding slides share:
% dac_bin_states, dac_bin_btran, dac_thermo_states and dac_thermo_tran.
%
% The originals draw plain circles with the code inside and straight arrows
% between them, with the two directions of a pair drawn side by side rather
% than as one double-headed arrow. \dacEdge offsets each arrow off the centre
% line so a pair reads as two arrows.

% #1 node name, #2 centre, #3 label. \dacStateWide takes a stacked label.
\newcommand{\dacState}[3]{%
  \node[draw, thick, circle, minimum size=1.5cm, inner sep=1pt] (#1) at #2 {#3};
}

\newcommand{\dacStateWide}[3]{%
  \node[draw, thick, circle, minimum size=1.9cm, inner sep=1pt,
        align=center] (#1) at #2 {#3};
}

% An arrow from one state to another. The two directions of a pair get their
% own lane by bending either side of the centre line — the originals draw
% them as parallel straight arrows, but drawn node to node the arrow is
% clipped at the circle border, which a hand-offset straight line is not.
\newcommand{\dacEdge}[2]{%
  \draw[thick, -{Latex[length=3mm]}] (#1) to[bend left=11] (#2);
}

% A single arrow between two states, on the centre line: used where the
% figure shows one transition rather than a pair.
\newcommand{\dacArrow}[2]{%
  \draw[thick, -{Latex[length=3mm]}] (#1) -- (#2);
}
